%----------------------------------------------------------------------------------------
% METADATA BLOCK (YAML)  — see inbox/README.md
%----------------------------------------------------------------------------------------
% uid: ""
% school: "lycee"
% level: "3eme"
% section: "Economie et Gestion"
% chapter: "Statistiques"
% title: "Achats dans une supérette"
% difficulty: "Facile"
% tags: ["statistiques", "classes", "classe-modale", "moyenne"]
% provenance: "own"
% batch: "2026-11-19-3eco-statistiques-n5"
%----------------------------------------------------------------------------------------

\newcommand{\R}{\mathbb{R}}
\newcommand{\vect}[1]{\overrightarrow{#1}}

% ---STATEMENT---
Montant des achats (en dinars) de $50$ clients d'une supérette :
\begin{center}
\begin{tabular}{|l|c|c|c|c|c|}\hline
Montant & $[0,20[$ & $[20,40[$ & $[40,60[$ & $[60,80[$ & $[80,100[$ \\\hline
Effectif & $6$ & $14$ & $20$ & $6$ & $4$ \\\hline
\end{tabular}\end{center}
\begin{enumerate}
  \item Donner la classe modale.
  \item Calculer les fréquences en pourcentage. Quel pourcentage des clients a dépensé moins de $40$ D ?
  \item Calculer le montant moyen des achats.
\end{enumerate}

% ---SOLUTION---
\begin{enumerate}[itemsep=14pt]
  \item Le plus grand effectif est $20$ : la classe modale est $[40,60[$.

  \item $f_i=\dfrac{n_i}{50}\times100=2\,n_i\,\%$ :
  \begin{center}
  \begin{tabular}{|l|c|c|c|c|c|}\hline
  Montant & $[0,20[$ & $[20,40[$ & $[40,60[$ & $[60,80[$ & $[80,100[$ \\\hline
  $f_i$ (\%) & $12$ & $28$ & $40$ & $12$ & $8$ \\\hline
  \end{tabular}\end{center}
  Moins de $40$ D : $12+28=40\,\%$ des clients.

  \item
  \begin{prop}[Série par classes]
  On remplace chaque classe $[a,b[$ par son \textbf{centre} \ib{$c_i=\dfrac{a+b}{2}$}, puis \ib{$\overline x=\dfrac{n_1c_1+\dots+n_pc_p}{N}$}.
  \end{prop}
  Centres : $10$, $30$, $50$, $70$, $90$.
  \[\overline x=\dfrac{6\times10+14\times30+20\times50+6\times70+4\times90}{50}\]
  \[\overline x=\dfrac{60+420+1\,000+420+360}{50}=\dfrac{2\,260}{50}=45{,}2\text{ D}\]
\end{enumerate}
